\documentclass[aspectratio=169, 12pt]{beamer}

% --- Language and Fonts ---
\usepackage{fontspec}
\usepackage[english]{babel}
\usepackage{amsmath, amssymb, bm}
% tikz
\usepackage{tikz}
\usetikzlibrary{shapes, arrows.meta, positioning, shadows, calc}

% --- Bibliography (BibTeX + natbib) ---
\usepackage{natbib}
\bibliographystyle{chicago} 

% =========================================
% ===== keybox ====================
% =========================================
\usepackage[most]{tcolorbox}
\newtcolorbox{keybox}[1]{
  enhanced,
  colback=headbg!4,
  colframe=headbg,
  coltitle=headfg,
  colbacktitle=headbg,
  title=#1,
  fonttitle=\bfseries,
  boxrule=0.8pt,
  arc=1.5mm,
  left=2mm,
  right=2mm,
  top=1mm,
  bottom=1mm,
  before skip=0.2cm,
  after skip=0.2cm
}

\usepackage{booktabs}

% =========================================
% ===== Custom Footline (Page Number) =====
% =========================================


\setbeamercolor{footline}{fg=headbg}
\setbeamerfont{footline}{size=\scriptsize}

\setbeamertemplate{footline}{%
  \hfill%
  \usebeamercolor[fg]{footline}%
  \usebeamerfont{footline}%
  \insertframenumber\,/\,\inserttotalframenumber\kern1em\vskip3pt%
}

\setbeamertemplate{navigation symbols}{}

\usetheme{default}
\usepackage{xcolor}

\definecolor{headbg}{RGB}{30, 60, 120}
\definecolor{headfg}{RGB}{255, 255, 255}
\definecolor{accent}{RGB}{200, 50, 50}

\setbeamercolor{headline}{bg=headbg, fg=headfg}
\setbeamercolor{section in head/foot}{bg=headbg, fg=headfg}
\setbeamercolor{subsection in head/foot}{bg=headbg!80, fg=headfg}
\setbeamercolor{frametitle}{bg=headbg, fg=headfg}

% Step indicator command
\newcommand{\stepindicator}[1]{%
  \hfill {\scriptsize \textit{\textcolor{headbg}{Step #1 / 6}}}%
}
\DeclareMathOperator*{\argmin}{arg\,min}

\usepackage{hyperref}


% --- Title ---
\title[]{A First Course in Causal Inference Ch~7}
\author[Kodai Minesaki]{Kodai Minesaki}
\institute[Faculty of Engineering, The University of Tokyo]{Faculty of Engineering, The University of Tokyo}
\date{\today}

\begin{document}

% --- Title Slide ---
\begin{frame}
\titlepage
\end{frame}


\begin{frame}

  {
    \large
    Ch.~7 Matched-Pairs Experiment
  }
  
\end{frame}

\begin{frame}{Ch.~7 Matched-Pairs Experiment}
  \begin{itemize}
    \item The matched-pairs experiment (MPE) is the most extreme version of the SRE with only one treated unit and one control unit within each stratum.
    \item It has its own estimation and inference strategy
  \end{itemize}
  
\end{frame}

\begin{frame}{7.1 Design of the experiment and potential outcomes}
  \begin{itemize}
    \item Consider an experiment with $2n$ units.
    \item If we have predictive covariates to the outcomes, we can pair units based on the similarity of covariates.
    \item With many covariates, we can define pairwise distances between units and the form pairs based on these distances.
    \item We assume that the pairs are formed based on the covariates and the discuss the subsequent design and analysis issues.
  \end{itemize}
\end{frame}

\begin{frame}{7.1 Design of the experiment and potential outcomes}
    Setup
    \begin{itemize}
        \item Let $(i,j)$ index the unit $j$ in pair $i$, where $i=1,\dots,n$ and $j=1,2$.
        \item Unit $(i,j)$ has potential outcomes $Y_{ij}(1)$ and $Y_{i,j}(0)$ under the treatment and control, respectively.
        \item Within each pair, we randomly assign one into the treatment group ($Z_i=1$) and the control group ($Z_i=0$).
    \end{itemize}
\end{frame}

\begin{frame}{7.1 Design of the experiment and potential outcomes}
    \small
    \begin{keybox}{Definition~7.1 (MPE)}
        \label{def-7-1}
        The treatment assignments satisfy
        \[
        (Z_i)_{i=1}^n \overset{\text{i.i.d.}}{\sim} \operatorname{Bernoulli}\left(\tfrac{1}{2}\right).
        \]
        The observed outcomes within pair $i$ are
        \[
        Y_{i1} = Z_i Y_{i1}(1) + (1 - Z_i) Y_{i1}(0) =
        \begin{cases}
            Y_{i1}(1), & \text{if } Z_i = 1, \\
            Y_{i1}(0), & \text{if } Z_i = 0,
        \end{cases}
        \]
        and
        \[
        Y_{i2} = Z_i Y_{i2}(0) + (1 - Z_i) Y_{i2}(1) =
        \begin{cases}
            Y_{i2}(0), & \text{if } Z_i = 1, \\
            Y_{i2}(1), & \text{if } Z_i = 0.
        \end{cases}
        \]
        So the observed data are $(Z_i,Y_{i1},Y_{i2})_{i=1}^n$
        \label{def-7-1}
    \end{keybox}
\end{frame}

\begin{frame}{7.3 Neymanian inference}
    \begin{itemize}
        \item The average causal effect within pair $i$ is 
        \[
        \tau_i = \frac{1}{2}\left[\{Y_{i1}(1)-Y_{i2}(0)\} + \{Y_{i2}(1)-Y_{i1}(0)\}\right].
        \]
        \item and the average causal effect for all units is
        \[
        \tau = \frac{1}{n}\sum_{i=1}^{n}\tau_i = \frac{1}{2n}\sum_{i=1}^{n}\sum_{j=1}^{2}\{Y_{ij}(1)-Y_{ij}(0)\}.
        \]
        \item Both $\hat{\tau}_i$ and $\hat{\tau}$ are unbiased for $\tau_i$ and $\tau$, respectively.
    \end{itemize}
\end{frame}

\begin{frame}{7.3 Neymanian inference}
    We can calculate the variance of $\hat{\tau}$.

    For simplicity, let $A_i=Y_{i1}(1)-Y_{i2}(0)$ and $B_i=Y_{i2}(1)-Y_{i1}(0)$. Then
    \[
    \hat{\tau}_i = Z_iA_i + (1-Z_i)B_i = Z_i(A_i-B_i)+B_i.
    \]
    Check the unbiasedness of $\hat{\tau}_i$: $$\mathbb{E}[\hat{\tau}_i] = \mathbb{E}[Z_i](A_i-B_i)+B_i = \dfrac{1}{2}(A_i-B_i)+B_i = \dfrac{A_i+B_i}{2}.\quad (\because Z_i\sim \text{Bernoulli}(1/2).)$$
    Here, $\operatorname{var}(\hat{\tau}_i) = \operatorname{var}(Z_i(A_i-B_i)+B_i) = \operatorname{var}(Z_i)(A_i-B_i)^2=\dfrac{1}{4}(A_i-B_i)^2.$ Therefore, 
    \[
    \operatorname{var}(\hat{\tau}) = \operatorname{var}\left(\frac{1}{n}\sum_{i=1}^{n}\operatorname{var}(\hat{\tau}_i)\right) = \frac{1}{4n^2}\sum_{i=1}^{n}\{Y_{i1}(1)-Y_{i2}(0) + Y_{i1}(0) - Y_{i2}(1)\}. \quad\Box
    \]
\end{frame}

\begin{frame}{7.3 Neymanian inference}
    \begin{itemize}
        \item We cannot follow the same strategy under the SRE to estimate the variance of $\hat{\tau}$ under the MPE.
        \item In MPE, the within-pair sample variances of the outcomes are not well defined.
        \begin{itemize}
            \item In SRE, we calculate variances for a treatment group and a control group in each stratum but in MPE, there is only one in both groups in each stratum.
        \end{itemize}
        \item Is it possible to estimate the variance of $\hat{\tau}$ in the MPE?
    \end{itemize}
\end{frame}

\begin{frame}{7.3 Neymanian inference}
    \begin{itemize}
        \item Think about the classic i.i.d. sampling:
        \item Suppose the $\hat{\tau}_i$'s are i.i.d. with mean $\mu$ and $\sigma^2$.
        \item Then the variance of $\hat{\tau}=n^{-1}\sum_{i=1}^{n}\hat{\tau}_i$ is $\sigma^2/n$.
        \item An unbiased estimator for $\sigma^2$ is $(n-1)^{-1}\sum_{i=1}^{n}(\hat{\tau}_i-\hat{\tau})^2$.  
        \item So the unbiased estimator for $\operatorname{var}(\hat{\tau})$ is 
        \begin{equation}
            \widehat{V}= \frac{1}{n(n-1)}\sum_{i=1}^{n}(\hat{\tau}_i-\hat{\tau})^2.
            \label{variance-estimator-mpe}
        \end{equation}

        \item We derive $\widehat{V}$ from a completely different setting from the MPE and a different motivation, so does it work for the MPE?
    \end{itemize}
\end{frame}

\begin{frame}{7.3 Neymanian inference}
    \begin{keybox}{Theorem~7.1}
        Under the MPE, $\widehat{V}$ defined in \eqref{variance-estimator-mpe} is a conservative estimator for the true variance of $\hat{\tau}$:
        \[
        \mathbb{E}[\widehat{V}] - \operatorname{var}(\hat{\tau}) = \frac{1}{n(n-1)}\sum_{i=1}^{n}(\tau_i-\tau)^2\geq 0.
        \]
        If the $\tau_i$'s are constant across pairs, then $\mathbb{E}[\widehat{V}] = \operatorname{var}(\hat{\tau})$.
        \label{thm-7-1}
    \end{keybox}
    \begin{itemize}
        \item Thm.~7.1 states that under the MPE,
        \begin{enumerate}
            \item $\widehat{V}$ is a conservative variance estimator in general;
            \item $\widehat{V}$ is unbiased if the average causal effecsts are constant across pairs.
        \end{enumerate}
    \end{itemize}
\end{frame}

\begin{frame}{7.3 Neymanian inference}
    Proof of Thm.~7.1:

    Recall $\sum_{i=1}^{n}(a_i-\bar{a})^2=\sum_{i=1}^{n}a_i^2-n\bar{a}^2$ and $\operatorname{var}[W]=\mathbb{E}[W^2]-(\mathbb{E}[W])^2$. Use them in the following manipulation:
    \begin{align*}
        n(n-1)\mathbb{E}[\widehat{V}] 
        &= \mathbb{E}\left[\sum_{i=1}^{n}(\hat{\tau}_i-\hat{\tau})^2\right]
        = \mathbb{E}\left[\sum_{i=1}^{n}\hat{\tau}_i^2-n\hat{\tau}^2\right]\\
        &= \sum_{i=1}^{n}\{\operatorname{var}(\hat{\tau})+\tau_i^2\} - n\{\operatorname{var}(\hat{\tau})+\tau^2\}\\
        &= n^2\operatorname{var}(\hat{\tau}) - n\operatorname{var}(\hat{\tau}) + \sum_{i=1}^{n}(\tau_i-\tau)^2.
    \end{align*}
    Therefore,
    \[
    \mathbb{E}[\widehat{V}] = \operatorname{var}(\hat{\tau}) + \{n(n-1)\}^{-1}\sum_{i=1}^{n}(\tau_i-\tau)^2\geq 0.\quad \Box
    \]
\end{frame}

\begin{frame}{7.3 Neymanian inference}
    Both the point estimator $\hat{\tau}$ and the variance estimator $\widehat{V}$ can be conveniently obtained by OLS:
    \begin{keybox}{Proposition 7.1}
        $\hat{\tau}$ and $\widehat{V}$ are identical to the coefficient and variance estimator of the intercept from the OLS fit of the vector $(\hat{\tau}_1,\hat{\tau}_2,\dots,\hat{\tau}_n)^T$ on the intercept only.
    \end{keybox}
\end{frame}

\begin{frame}{7.3 Neymanian inference}

\end{frame}

% =========================

% =========================
\section{}
\begin{frame}[allowframebreaks, noframenumbering]{References}
  \beamerdefaultoverlayspecification{}
  \small
  \bibliography{../references}
\end{frame}

\end{document}